% ============================================================
% APPENDIX A: TABLES (CSV files written by R/Q*.R in tables/)
% ============================================================
\section{Tables}\label{app:tables}

\subsection{Question 1}

% TODO may be rm
\csvtable{q1a_variables}
  {a) Variables of the data set.}
  {tab:q1a-variables}

% TODO may rm and just output log since boxplot
\csvtable{q1b_summary}
  {b) Summary statistics of the numerical variables (original data).}
  {tab:q1b-summary}

\csvtable{q1b_outliers}
  {b) Outliers ($1.5 \times \mathrm{IQR}$ rule) and zero values of the numerical variables.}
  {tab:q1b-outliers}

  % TODO drop the BMI, keep only variable from dataset

  % BMI -> height
\csvtable[\footnotesize]{q1b_removed}
  {b) Observations removed as impossible values or outside the adult reference ranges of \autoref{tab:q1b-reference} (blood pressure values as \texttt{systolic}/\texttt{diastolic}; only their number if more than 10).}
  {tab:q1b-removed}

% TODO may rm and just output log since boxplot
\csvtable{q1b_summary_clean}
  {b) Summary statistics of the numerical variables (cleaned data).}
  {tab:q1b-summary-clean}

\begin{table}[htbp]
  \centering
  \footnotesize
  \caption{b) Reference ranges of the health measures for adults. Sources: BMI \citep{clevelandclinic_bmi}, \texttt{systolic} and \texttt{diastolic} \citep{clevelandclinic_bp}, and \texttt{body\_fat\_pct} \citep{ace_bodyfat}.}
  \label{tab:q1b-reference}
  \begin{tabular}{lll}
    \toprule
    Measure & Category & Range \\
    \midrule
    BMI (kg/m$^2$), from
      & Underweight           & $< 18.5$ \\
    \quad \texttt{height\_cm}, \texttt{weight\_kg}
      & Normal weight         & 18.5--24.9 \\
      & Overweight            & 25--29.9 \\
      & Obesity               & $\geq 30$ \\
      & Obesity class III     & $\geq 40$ \\
    \midrule
    \texttt{body\_fat\_pct} (\%),
      & Essential fat         & 2--5 / 10--13 \\
    \quad men / women
      & Athletes              & 6--13 / 14--20 \\
      & Fitness               & 14--17 / 21--24 \\
      & Average               & 18--24 / 25--31 \\
      & Obesity               & $\geq 25$ / $\geq 32$ \\
    \midrule
    \texttt{systolic} / \texttt{diastolic}
      & Low (hypotension)     & $< 90$ / $< 60$ \\
    \quad (mmHg)
      & Normal                & $< 120$ and $< 80$ \\
      & Elevated              & 120--129 and $< 80$ \\
      & Stage 1 hypertension  & 130--139 or 80--89 \\
      & Stage 2 hypertension  & $\geq 140$ or $\geq 90$ \\
      & Hypertensive crisis   & $> 180$ and/or $> 120$ \\
    \bottomrule
  \end{tabular}
\end{table}

\csvtable{q1b_shapiro}
  {b) Shapiro-Wilk test of the numerical variables (random sample of 5,000 observations of the cleaned data).}
  {tab:q1b-shapiro}

\csvtable{q1c_summary}
  {c) Descriptive statistics of systolic and diastolic blood pressure by gender (mmHg).}
  {tab:q1c-summary}

\csvtable{q1d_t_test}
  {d) One-sample t-tests ($H_0$: mean of 120\,mmHg) and 95\% confidence intervals for the mean systolic blood pressure by gender (mmHg).}
  {tab:q1d-t-test}

% TODO may rm and just output log since bar plot 
\csvtable{q1e_age_groups}
  {e) Frequency of each age group (\texttt{age\_cat}).}
  {tab:q1e-age-groups}
% TODO may rm and just output log since bar plot 
\csvtable{q1f_age_class}
  {f) Contingency table of age group and performance class.}
  {tab:q1f-age-class}
% TODO may rm and just output log since bar plot 
\csvtable{q1f_age_gender}
  {f) Contingency table of age group and gender.}
  {tab:q1f-age-gender}

\csvtable{q1f_chisq}
  {f) Chi-square tests of independence.}
  {tab:q1f-chisq}

\csvtable{q1g_shapiro}
  {g) Shapiro-Wilk normality tests of the physical performance variables by gender (random sample of 5,000 males; all females).}
  {tab:q1g-shapiro}

\csvtable{q1g_f_test}
  {g) F-tests for equality of variances and t-test chosen.}
  {tab:q1g-f-test}

\csvtable{q1g_t_test}
  {g) Two-sample t-test results for the physical performance variables by gender, and 95\% confidence intervals of the difference of means (males minus females).}
  {tab:q1g-t-test}
  
\clearpage
\subsection{Question 2}

\csvtable[\footnotesize]{q2a_summary_age}
  {a) Summary statistics of the four variables in each age group (\texttt{age\_cat}).}
  {tab:q2a-summary-age}

\csvtable[\footnotesize]{q2a_summary_class}
  {a) Summary statistics of the four variables in each performance class (\texttt{class}).}
  {tab:q2a-summary-class}

\csvtable{q2b_anova}
  {b) One-way ANOVA of the four variables by \texttt{age\_cat} and by \texttt{class} ($\mathrm{Eta\_sq} = \eta^2 = \mathrm{SSB}/\mathrm{SST}$).}
  {tab:q2b-anova}

\csvtable[\footnotesize]{q2b_assumptions}
  {b) Assumptions of the eight one-way ANOVA models (Shapiro-Wilk on a random sample of 5,000 residuals; \texttt{Var\_ratio}: largest over smallest group variance).}
  {tab:q2b-assumptions}

\csvtable[\scriptsize]{q2b_posthoc_age}
  {b) Post-hoc comparisons of each pair of age groups (\texttt{age\_cat}).}
  {tab:q2b-posthoc-age}

\csvtable[\scriptsize]{q2b_posthoc_class}
  {b) Post-hoc comparisons of each pair of performance classes (\texttt{class}).}
  {tab:q2b-posthoc-class}

\csvtable{q2b_posthoc_summary}
  {b) Number of the 6 pairs of groups that differ at a 95\% confidence level ($p \leq 0.05$) according to each post-hoc test.}
  {tab:q2b-posthoc-summary}

\csvtable{q2c_anova}
  {c) Two-way ANOVA (Type~II tests) of \texttt{sit\_ups} and \texttt{broad\_jump\_cm} by \texttt{age\_cat} and \texttt{gender}.}
  {tab:q2c-anova}

\csvtable[\footnotesize]{q2c_assumptions}
  {c) Assumptions of the two two-way ANOVA models across the eight \texttt{age\_cat} $\times$ \texttt{gender} groups.}
  {tab:q2c-assumptions}

\csvtable[\scriptsize]{q2c_posthoc}
  {c) Pairwise post-hoc comparisons of \texttt{age\_cat} within each level of \texttt{gender} for \texttt{sit\_ups} and \texttt{broad\_jump\_cm}.}
  {tab:q2c-posthoc}

\clearpage
\subsection{Question 3}

\csvtable{q3a_correlation}
  {a) Correlation matrix of \texttt{broad\_jump\_cm} and the numerical predictors.}
  {tab:q3a-correlation}

\csvtable{q3a_slr}
  {a) Simple linear regression models of \texttt{broad\_jump\_cm}.}
  {tab:q3a-slr}

\csvtable[\scriptsize]{q3a_assumptions}
  {a) Assumptions of the best simple model, without and with the Box-Cox transformation of the response (Shapiro-Wilk on a random sample of 5,000 residuals; outliers: $|$standardized residual$| > 2$).}
  {tab:q3a-assumptions}

\csvtable[\scriptsize]{q3b_selection}
  {b) Models of \texttt{broad\_jump\_cm} with two of the three numerical predictors; partial F-test against the model with the three predictors.}
  {tab:q3b-selection}

\csvtable[\scriptsize]{q3b_comparison}
  {b) Simple model of a) and multiple model of b), both without transformation; partial F-test of the predictor added in b).}
  {tab:q3b-comparison}

\csvtable[\scriptsize]{q3b_assumptions}
  {b) Assumptions of the multiple linear regression model (Shapiro-Wilk on a random sample of 5,000 residuals; outliers: $|$standardized residual$| > 2$; largest VIF of the predictors).}
  {tab:q3b-assumptions}

\csvtable{q3c_age_coef}
  {c) i) Coefficients of the model of b) with \texttt{age\_cat} (reference category \texttt{<30}) and 95\% confidence intervals.}
  {tab:q3c-age-coef}

\csvtable{q3c_gender_coef}
  {c) ii) Coefficients of the model of b) with \texttt{gender} (reference category \texttt{Female}) and 95\% confidence intervals.}
  {tab:q3c-gender-coef}

\csvtable[\scriptsize]{q3c_summary}
  {c) Overall summary of the two models of c).}
  {tab:q3c-summary}

\csvtable[\scriptsize]{q3e_validity}
  {e) Validity checks of the final model (model of b) with \texttt{gender}); Shapiro-Wilk on a random sample of 5,000 residuals, outliers with $|$standardized residual$| > 2$, SD: standard deviation of the residuals of women and of men.}
  {tab:q3e-validity}

% Extension: heteroscedasticity-consistent (HC3) standard errors, not in the slides
\csvtable{q3e_robust}
  {e) Coefficients of the final model with heteroscedasticity-consistent (HC3) standard errors, which do not assume homoscedasticity.}
  {tab:q3e-robust}

% ============================================================
% APPENDIX B: FIGURES (PNG files written by R/Q*.R in figures/)
% ============================================================
\clearpage
\section{Figures}\label{app:figures}

\subsection{Question 1}

\begin{figure}[htbp]
  \centering
  \includegraphics[width=\linewidth]{q1b_boxplots}
  \caption{b) Boxplots of the numerical variables (original data).}
  \label{fig:q1b-boxplots}
\end{figure}

\begin{figure}[htbp]
  \centering
  \includegraphics[width=\linewidth]{q1b_boxplots_clean}
  \caption{b) Boxplots of the numerical variables (cleaned data).}
  \label{fig:q1b-boxplots-clean}
\end{figure}

\begin{figure}[htbp]
  \centering
  \includegraphics[width=\linewidth]{q1b_histograms}
  \caption{b) Histograms of the numerical variables (cleaned data).}
  \label{fig:q1b-histograms}
\end{figure}

\begin{figure}[htbp]
  \centering
  \includegraphics[width=\linewidth]{q1b_qqplots}
  \caption{b) Normal Q-Q plots of the numerical variables (cleaned data; red line: normal distribution).}
  \label{fig:q1b-qqplots}
\end{figure}

\begin{figure}[htbp]
  \centering
  \includegraphics[width=\linewidth]{q1c_histograms}
  \caption{c) Histograms of systolic and diastolic blood pressure by gender.}
  \label{fig:q1c-histograms}
\end{figure}

\begin{figure}[htbp]
  \centering
  \includegraphics[width=\linewidth]{q1c_boxplots}
  \caption{c) Boxplots of systolic and diastolic blood pressure by gender.}
  \label{fig:q1c-boxplots}
\end{figure}

\begin{figure}[htbp]
  \centering
  \includegraphics[width=0.6\linewidth]{q1e_age_groups}
  \caption{e) Number of individuals in each age group.}
  \label{fig:q1e-age-groups}
\end{figure}

\begin{figure}[htbp]
  \centering
  \includegraphics[width=\linewidth]{q1f_proportions}
  \caption{f) Proportion of each performance class (left) and gender (right) in each age group.}
  \label{fig:q1f-proportions}
\end{figure}

\begin{figure}[htbp]
  \centering
  \includegraphics[width=\linewidth]{q1g_distributions}
  \caption{g) Density plots of the physical performance variables by gender.}
  \label{fig:q1g-distributions}
\end{figure}

\begin{figure}[htbp]
  \centering
  \includegraphics[width=\linewidth]{q1g_boxplots}
  \caption{g) Boxplots of the physical performance variables by gender.}
  \label{fig:q1g-boxplots}
\end{figure}
\clearpage
\subsection{Question 2}

\begin{figure}[htbp]
  \centering
  \includegraphics[width=\linewidth]{q2a_boxplots_age}
  \caption{a) Boxplots of the four variables by age group (\texttt{age\_cat}).}
  \label{fig:q2a-boxplots-age}
\end{figure}

\begin{figure}[htbp]
  \centering
  \includegraphics[width=\linewidth]{q2a_boxplots_class}
  \caption{a) Boxplots of the four variables by performance class (\texttt{class}).}
  \label{fig:q2a-boxplots-class}
\end{figure}

\begin{figure}[htbp]
  \centering
  \includegraphics[width=\linewidth]{q2b_qqplots}
  \caption{b) Normal Q-Q plots of the residuals of the eight one-way ANOVA models (top: \texttt{age\_cat}, bottom: \texttt{class}).}
  \label{fig:q2b-qqplots}
\end{figure}

\begin{figure}[htbp]
  \centering
  \includegraphics[width=\linewidth]{q2c_interaction}
  \caption{c) Interaction plots of the mean of \texttt{sit\_ups} (left) and \texttt{broad\_jump\_cm} (right) by \texttt{age\_cat} and \texttt{gender}.}
  \label{fig:q2c-interaction}
\end{figure}

\begin{figure}[htbp]
  \centering
  \includegraphics[width=\linewidth]{q2c_qqplots}
  \caption{c) Normal Q-Q plots of the residuals of the two two-way ANOVA models.}
  \label{fig:q2c-qqplots}
\end{figure}

\clearpage
\subsection{Question 3}

\begin{figure}[htbp]
  \centering
  \includegraphics[width=\linewidth]{q3a_scatter}
  \caption{a) \texttt{broad\_jump\_cm} against each numerical predictor, with the fitted simple regression line.}
  \label{fig:q3a-scatter}
\end{figure}

\begin{figure}[htbp]
  \centering
  \includegraphics[width=0.6\linewidth]{q3a_boxcox}
  \caption{a) Box-Cox profile log-likelihood of the best simple model, with the 95\% confidence interval of $\lambda$.}
  \label{fig:q3a-boxcox}
\end{figure}

\begin{figure}[htbp]
  \centering
  \includegraphics[width=\linewidth]{q3a_diagnostics}
  \caption{a) Residuals against fitted values (left) and normal Q-Q plot of the residuals (right) of the best simple model.}
  \label{fig:q3a-diagnostics}
\end{figure}

\begin{figure}[htbp]
  \centering
  \includegraphics[width=\linewidth]{q3b_diagnostics}
  \caption{b) Residuals against fitted values (left) and normal Q-Q plot of the residuals (right) of the multiple linear regression model.}
  \label{fig:q3b-diagnostics}
\end{figure}

\begin{figure}[htbp]
  \centering
  \includegraphics[width=0.8\linewidth]{q3b_pairs}
  \caption{b) Scatter plot matrix of \texttt{broad\_jump\_cm}, \texttt{grip\_force} and \texttt{sit\_ups}: scatter plots with the fitted line below the diagonal, correlations above it and histograms on the diagonal (whole-number values slightly jittered).}
  \label{fig:q3b-pairs}
\end{figure}

\begin{figure}[htbp]
  \centering
  \includegraphics[width=\linewidth]{q3b_vs_a}
  \caption{b) Observed against predicted \texttt{broad\_jump\_cm} of the simple model of a) and the multiple model of b), both without transformation; grey line: $y = x$.}
  \label{fig:q3b-vs-a}
\end{figure}

\begin{figure}[htbp]
  \centering
  \includegraphics[width=\linewidth]{q3c_categories}
  \caption{c) Predicted \texttt{broad\_jump\_cm} for each age group (top, model c i) and each gender (bottom, model c ii), against \texttt{grip\_force} with \texttt{sit\_ups} at its mean (left) and against \texttt{sit\_ups} with \texttt{grip\_force} at its mean (right).}
  \label{fig:q3c-categories}
\end{figure}

\begin{figure}[htbp]
  \centering
  \includegraphics[width=\linewidth]{q3d_fit}
  \caption{d) Observed against predicted \texttt{broad\_jump\_cm} of the model of b) and the two models of c); grey line: $y = x$.}
  \label{fig:q3d-fit}
\end{figure}

\begin{figure}[htbp]
  \centering
  \includegraphics[width=\linewidth]{q3e_residuals}
  \caption{e) Residual analysis of the final model: residuals against fitted values, normal Q-Q plot, histogram of the residuals with the normal density (blue) and residuals by gender.}
  \label{fig:q3e-residuals}
\end{figure}

% ============================================================
% APPENDIX C: R CODE AND OUTPUTS
% ============================================================
\clearpage
\section{R Code and Outputs}\label{app:code}

The analysis is done in R, with one script per question in the \texttt{R/} folder of the project: \texttt{R/Q1.R}, \texttt{R/Q2.R} and \texttt{R/Q3.R}, with one section per part of the question.

Each script writes three kinds of output: the tables as CSV files (Appendix~\ref{app:tables}), the charts as PNG files (Appendix~\ref{app:figures}), and its full console output in a log, \texttt{R/logs/<script>.log}.
